\documentclass[11pt,a4paper]{article}

\usepackage[T1]{fontenc}
\usepackage{lmodern}
\usepackage{amsmath,amssymb}
\usepackage[margin=2.4cm]{geometry}
\usepackage{booktabs}
\usepackage{array}
\usepackage[colorlinks=true,linkcolor=blue,urlcolor=blue]{hyperref}

\newcommand{\Mqq}{M_{q\bar q}^2}
\newcommand{\dd}{\mathrm{d}}
\newcommand{\xP}{x_{\mathbb{P}}}
\newcommand{\code}[1]{\texttt{#1}}
\newcommand{\opt}[1]{\texttt{-{}-#1}}

\title{Numerical evaluation of the DIS structure functions\\in the C++ code}
\author{}
\date{\today}

\begin{document}
\maketitle

\begin{abstract}
  This note describes how the C++ code in \code{cpp/} evaluates the
  leading-order inclusive DIS structure functions with a BK-evolved dipole: the
  optical-theorem result (OT) and the finite-energy (FE) constrained results
  with the dipole evolved to $x_B$ (4D) or to $\xP$ (5D). It explains the
  integration variables, the three ways the $\Mqq$ integral of the 5D cross
  section can be done, and the choice of integration routine, and it compares
  the results with the Python reference implementation in
  \code{src/small\_x\_physics}. The C++ code computes the same integrals as
  the Python code. With the default settings it agrees with the Python results
  within their statistical errors, and it reaches comparable precision 30--80
  times faster on 8 cores.
\end{abstract}

\tableofcontents

\section{The quantities}
\label{sec:quantities}

The code computes, for one quark flavour of mass $m$ and charge $Z_f$, the
structure functions
\begin{equation}
  F_\lambda = \frac{Q^2}{4\pi^2 \alpha_{\text{em}}} \, \sigma_\lambda ,
  \qquad \lambda = L, T, \qquad F_2 = F_L + F_T ,
\end{equation}
in three ways, exactly as \code{test\_bk\_integration.py}. With
$r^2 = |\mathbf{u} - \mathbf{u}'|^2 = u^2 + u'^2 - 2uu'\cos\theta$,
$W^2 = Q^2(1/x_B - 1)$ and $\epsilon^2 = Q^2 z(1-z) + m^2$:

\paragraph{OT} (\code{BK\_OT\_CS.py}), with the dipole at $Y = \ln(x_0/x_B)$:
\begin{equation}
  \sigma_\lambda^{\text{OT}} = \int_{r_{\min}}^{r_{\max}} \dd r
    \int_{z_{\min}}^{1 - z_{\min}} \dd z \;
    \frac{\sigma_0}{2} \, \frac{1}{4\pi} \, \frac{2\pi r}{z(1-z)} \,
    |\psi_\lambda(r, z)|^2 \; 2\big[1 - S_Y(r)\big] .
\end{equation}

\paragraph{FE 4D} (\code{BK\_FE\_4D\_evolve\_in\_xB.py}), with the dipole at
$Y = \ln(x_0/x_B)$:
\begin{equation}
  \sigma_\lambda^{\text{4D}}
  = \int \dd u \, \dd u' \, \dd\theta \, \dd z \;
    P_\lambda(u, u', \theta, z) \; T_Y(u, u', \theta, z) \;
    \frac{\zeta J_1(\zeta)}{2\pi r^2} ,
  \qquad
  \zeta = r \sqrt{z(1-z) W^2 - m^2} .
  \label{eq:fe4d}
\end{equation}

\paragraph{FE 5D} (\code{BK\_FE\_5D\_evolve\_in\_xP.py}), with the dipole at
$Y(\Mqq) = \ln(x_0/\xP)$ and $\xP = (\Mqq + Q^2)/(W^2 + Q^2)$:
\begin{equation}
  \sigma_\lambda^{\text{5D}}
  = \int \dd u \, \dd u' \, \dd\theta \, \dd z
    \int_{4m^2}^{W^2} \dd\Mqq \;
    P_\lambda \; T_{Y(\Mqq)} \;
    \Theta\big(k^2\big) \, \frac{z(1-z)}{4\pi} \, J_0(k r) ,
  \qquad
  k^2 = z(1-z)\Mqq - m^2 .
  \label{eq:fe5d}
\end{equation}

Here $P_\lambda = \frac{\sigma_0}{2} \frac{1}{4\pi}
\frac{2\pi u u'}{z(1-z)} |\psi_\lambda(u, u', \theta, z)|^2$ collects the
light-cone wavefunction and its conjugate, and
\begin{equation}
  T_Y = 1 - S_Y(u) - S_Y(u') + S_{4,Y}(u, u', \theta, z)
\end{equation}
is the target amplitude, built from the BK dipole $S_Y$ and the Gaussian
quadrupole $S_{4,Y}$ at finite or large $N_c$. The limits are those of the
Python script: $r, u, u' \in [10^{-6}, 10]\ \text{GeV}^{-1}$,
$z \in [10^{-8}, 1 - 10^{-8}]$, $\theta \in [0, 2\pi]$.

All numbers in this note are for $Q^2 = 2\ \text{GeV}^2$, $x_B = 10^{-3}$,
$m = 0.14\ \text{GeV}$, $Z_f^2 = 2/3$, $\sigma_0 = 2 \cdot 2.57 \cdot 18.81\
\text{GeV}^{-2}$, $x_0 = 0.01$, the BK solution with MV initial condition in
\code{data/mv.dat}, and the finite-$N_c$ quadrupole unless stated otherwise.

\section{Building blocks}

All of these agree with the Python implementation point by point. The unit tests
compare them at about 2000 points each, generated by
\code{tests/reference/generate\_reference.py}.

\begin{description}
  \item[BK dipole.] The Python code uses scipy's \code{RectBivariateSpline}
    with $s = 0$. That is the unique bicubic tensor-product interpolant with
    not-a-knot end conditions, so the C++ code rebuilds it from the data by 1D
    not-a-knot splines along each axis, stored as one polynomial per grid cell.
    Arguments outside the grid are clamped, $N$ is clipped to $[0, 1]$ and
    $S = \max(1 - N, 10^{-14})$, as in Python. Agreement: $10^{-13}$ in $N$.
    The interpolation is a template parameter, and a bilinear alternative is
    provided.
  \item[Large-$x$ extrapolation.] Optionally (\opt{large-x-power} $p$), the
    dipole at $x > x_0$ is
    $N(r, x) = N(r, x_0)\,[(1-x)/(1-x_0)]^p$ instead of the Python behaviour
    of keeping it at $x_0$. Only FE 5D evaluates the dipole at $x > x_0$. There
    is no Python counterpart.
  \item[Wavefunctions.] The Bessel functions $K_0$, $K_1$, $J_0$, $J_1$ come
    from Boost.Math, or from GSL with \code{-Duse\_boost=OFF}. Agreement:
    $10^{-12}$ relative.
  \item[Quadrupole.] Dominguez et al.\ eq.~B21 from the six pair dipoles,
    including the $\Delta = 0$ branch and the same $10^{-14}$ floor in the
    logarithms. Tested with an analytic GBW dipole and with the BK dipole.
    Agreement: $10^{-10}$.
  \item[Integrands.] The OT, FE 4D and FE 5D integrands agree to $10^{-9}$,
    relative to the integrand's scale.
\end{description}

L and T share everything except the wavefunction, so each integrand returns
both, and one integration computes them as two components. Computing them in
separate integrations, as Python does, gave 10--15\% smaller FE errors at equal
points per polarization but cost twice as much, so at equal time the combined
integration is better.

\section{Integration variables}
\label{sec:variables}

The C++ code integrates over the same domains as Python, but by default in
different variables. \opt{variables python} restores the Python variables.

\paragraph{Relative coordinates (FE).}
The phase-space kernel of \eqref{eq:fe4d} is sharply peaked at
$\mathbf{u} = \mathbf{u}'$, with a width of about $1/\sqrt{z(1-z)W^2}$, roughly
$0.05\ \text{GeV}^{-1}$ here. Away from the peak it falls only like
$r^{-3/2}$, oscillating. Its integral converges, but its square does not:
\emph{in the Python variables $(u, u', \theta)$ the integrand has infinite
variance} near the diagonal $u = u'$, $\theta = 0$. That is a diagonal of the
integration box, which an adaptive Monte Carlo cannot follow. Section~\ref{sec:python}
shows the consequences: estimates drift upwards with the number of points, and
the errors vegas reports are too small.

The C++ code integrates instead over $u$, the separation
$r = |\mathbf{u} - \mathbf{u}'|$ on a logarithmic scale, and the angle $\phi$
of $\mathbf{u} - \mathbf{u}'$ relative to $\mathbf{u}$. Since
$\dd^2u\,\dd^2u' = \dd^2u\,\dd^2r$, the integration measure changes as
\begin{equation}
  \dd u \, \dd u' \, \dd\theta \; u u' = \dd u \, \dd r \, \dd\phi \; u r ,
\end{equation}
so the integrand gains a factor $r/u'$. The original domain
$u' \in [r_{\min}, r_{\max}]$ is kept exactly by setting the integrand to zero
outside it. The disc $r < 10^{-10}\ \text{GeV}^{-1}$ is left out of the
logarithmic scale; it changes the result by a relative
$\sim (10^{-10}\,\text{GeV}^{-1} \cdot M)^2 \sim 10^{-15}$. $r^2$ is computed
directly from $r$, not from $u^2 + u'^2 - 2uu'\cos\theta$, which would cancel
catastrophically at $r \ll u$.

\paragraph{Symmetries.}
All integrands are invariant under $z \to 1 - z$, which swaps quark and
antiquark, and under $\theta \to -\theta$, a reflection. So $z$ runs over
$[z_{\min}, 1/2]$ and the angle over $[0, \pi]$, and the result is multiplied
by 4 (by 2 for OT).

\paragraph{Logarithmic sizes.} $r$ (OT) and $u$ (FE) are sampled on a
logarithmic scale.

\medskip
With Vegas at $10^6$ points, these variables make the FE 4D errors 1.8--1.9
times smaller than in the Python variables. Over repeated seeds the reported
errors match the observed scatter (ratio 0.6--1.5 over 8 seeds). For OT the
variables make no difference with Vegas; see Section~\ref{sec:integrators}.

\section{The \texorpdfstring{$\Mqq$}{M2} integral of FE 5D}
\label{sec:fe5d}

At fixed $u, u', \theta, z$, only $T$ depends on $\Mqq$ in \eqref{eq:fe5d},
through $Y(\Mqq)$. With $s = k^2$, $\dd s = z(1-z)\,\dd\Mqq$, and
$K^2 = z(1-z)W^2 - m^2$, the inner integral is
\begin{equation}
  I = \frac{1}{4\pi} \int_0^{K^2} \dd s \; J_0\big(r\sqrt{s}\big) \, T\big(Y(s)\big) .
  \label{eq:inner}
\end{equation}
$J_0$ oscillates over the whole range (up to $K r \sim 400$ here), and the
narrow peak of the 4D kernel builds up only by cancellation. That makes the
original form hard for any integrator. Three forms are available
(\opt{fe5d-method}).

\subsection{Original form (\code{original})}
The 5D integral \eqref{eq:fe5d} as written, as in Python. It is used with
\opt{variables python} and for cross-checks.

\subsection{Integration by parts (\code{by-parts})}
Using $\frac{\dd}{\dd s}\big[2\sqrt{s}\,J_1(r\sqrt{s})/r\big] = J_0(r\sqrt{s})$,
\begin{equation}
  I = \frac{K J_1(K r)}{2\pi r}\, T(Y_W)
    + \frac{1}{4\pi} \int \dd\Mqq \; \frac{2kJ_1(kr)}{r} \,
      \frac{1}{\Mqq + Q^2} \, \frac{\dd T}{\dd Y} ,
  \qquad Y_W = Y(W^2) = \ln x_0 .
\end{equation}
The first term is the FE 4D integrand \eqref{eq:fe4d} with the dipole at
$\xP = 1$. The second, the correction, vanishes wherever the dipole does not
depend on $Y$. For the Python behaviour (dipole kept at $x_0$ for $x > x_0$),
that is all of $\xP > x_0$, so the correction only covers
$\Mqq \lesssim 18\ \text{GeV}^2$ instead of up to $W^2 \approx 2000\
\text{GeV}^2$. There it barely oscillates. $\dd T/\dd Y$ is a central finite
difference with step $10^{-4}$, kept inside the dipole's rapidity range. The
derivation is in \code{doc/fe5d\_by\_parts.pdf}. With Vegas this reduces the
FE 5D error 12--16 times at equal points.

With the large-$x$ extrapolation the dipole depends on $Y$ all the way to
$\xP = 1$. Then the boundary term vanishes ($N = 0$ at $\xP = 1$, so $T = 0$),
and the correction spans the whole range and oscillates again. In that case
Cuhre's error estimate fails on it: results were off by about 20 times the
estimate. The code warns if this combination is requested.

\subsection{Filon rule (\code{filon})}
The inner integral \eqref{eq:inner} is computed at each point
$(u, u', \theta, z)$, so that the remaining integral is a smooth 4D one.
\begin{enumerate}
  \item The rapidity range $[Y(K^2), Y(0)]$ is split at the dipole's
    \emph{breakpoints}, where $\dd T/\dd Y$ may jump: the ends of the tabulated
    range and, with the extrapolation, $Y = 0$. On a piece where the dipole does
    not depend on $Y$, $T$ is constant and the integral of $J_0$ is exact.
  \item On the other pieces, $T$ is evaluated at 8 Chebyshev--Lobatto nodes and
    interpolated with the barycentric formula. The nodes are in $Y$ for
    $\xP < x_0$, where BK evolution is smooth in $Y$, and in $\xP$ for
    $\xP > x_0$, where the extrapolation is a smooth function of $x$. In $Y$ it
    would be stretched exponentially: the first version, with nodes in $Y$, was
    off by up to $10^{-3}$ of the inner integral in typical configurations, and
    by 5\% where the integral nearly cancels.
  \item Each piece is divided into 8 panels. On each panel the interpolant is
    replaced by a cubic in $s$ through four points, and $J_0(r\sqrt{s})$ times
    the cubic is integrated exactly. This needs the moments
    $\int s^n J_0(r\sqrt{s})\,\dd s$, $n \le 3$, whose antiderivative is, with
    $x = r\sqrt{s}$,
    \begin{equation}
      A_n(s) = s^{n+1} \sum_{j\ge0} \frac{(-x^2/4)^j}{(j!)^2 (n+j+1)}
      \quad (x < 4), \qquad
      A_n(s) = \frac{2\, I_{2n+1}(x)}{r^{2n+2}} \quad (x \ge 4),
    \end{equation}
    where $I_m(x) = \int_0^x y^m J_0(y)\,\dd y$ follows from $I_1 = xJ_1$ and
    $I_m = x^m J_1 + (m-1)x^{m-1}J_0 - (m-1)^2 I_{m-2}$. The series avoids the
    cancellation of the recurrence at small $x$.
\end{enumerate}
The accuracy is that of interpolating $T$, independently of how much $J_0$
oscillates. A unit test compares the Filon integrand with a brute-force 1D
quadrature of \eqref{eq:inner} for the frozen and the extrapolated dipole.
Each point costs about 30 FE 4D evaluations.

\subsection{Automatic choice (\code{auto}, the default)}
\code{by-parts} if the dipole stops depending on $Y$ inside the $\Mqq$ range
(the Python behaviour), \code{filon} if it depends on $Y$ up to $\xP = 1$ (the
extrapolation). For a frozen dipole the Filon result keeps the sharp
$r \to 0$ peak created by $T$ jumping to a nonzero constant at $W^2$, so Cuhre
needs millions of the expensive points (12.8~s against 1.5~s by parts). With
the extrapolation, $T$ vanishes smoothly at $\xP = 1$, and Cuhre converges in
about $10^5$ points.

\section{Integration routines}
\label{sec:integrators}

\opt{integrator} selects \textbf{Cuhre} (the default) or \textbf{Vegas} from
Cuba 4.2.2.
\begin{description}
  \item[Cuhre] is deterministic adaptive cubature. It subdivides the region
    until the estimated error is below \opt{epsrel} (default $10^{-3}$) times the
    result. Its error estimate compares cubature rules of different degree. It
    is not statistical, so it cannot be checked by changing the seed; it was
    checked by repeating with a ten times smaller \opt{epsrel}, and the change
    was always within the estimate. Cuhre keeps all subregions in memory and
    slows down beyond a few $10^7$ evaluations.
  \item[Vegas] is adaptive Monte Carlo, run like Python's vegas: 10 warm-up
    iterations of $N/10$ points that train the grid and are discarded, then 20
    iterations of $N$ points (\opt{mcpoints} $N$). Cuba implements classic
    Vegas, without the stratified sampling of Python's vegas package.
\end{description}

Table~\ref{tab:integrators} compares the Cuba routines at the same budget of
$2.1\times10^7$ evaluations, the budget of Vegas at $N = 10^6$, with the
default variables and the frozen dipole.

\begin{table}[ht]
  \centering
  \footnotesize
  \setlength{\tabcolsep}{4pt}
  \begin{tabular}{lllll}
    \toprule
    & OT & FE 4D & FE 5D correction & time vs.\ Vegas \\
    \midrule
    Vegas                  & $1.6\times10^{-5}$ & $1.8\times10^{-4}$ & $5.9\times10^{-6}$ & 1 \\
    Suave, $n_{\text{new}} = 10^4$  & biased ($6\sigma$) & $1.6\times10^{-4}$ & biased ($5$--$8\sigma$) & $\approx 1.3$ \\
    Suave, $n_{\text{new}} = 10^5$  & $2.7\times10^{-6}$ & $1.7\times10^{-4}$ & $4.8\times10^{-6}$ & $\approx 1.3$ \\
    Divonne                & $5.6\times10^{-8}$, FT off $3\sigma$ & $8.0\times10^{-5}$ & --- & 25--450 \\
    Cuhre                  & converged & $9.5\times10^{-6}$ & $2.4\times10^{-7}$ & $\approx 1$ \\
    \bottomrule
  \end{tabular}
  \caption{Errors of $F_L$ from the Cuba routines at $2.1\times10^7$
    evaluations. ``Biased'' means off
    from the reference by the stated multiple of the routine's own error.}
  \label{tab:integrators}
\end{table}

\paragraph{OT.} Classic Vegas cannot match the stratified sampling of Python's
vegas for this smooth 2D integral: at equal points its error is 50--70 times
larger, and no change of variables helped. Cuhre gives the integral to
$10^{-9}$ in 0.1~s.

\paragraph{FE 5D with the extrapolation.} This is the case that matters most
for future work. Table~\ref{tab:suave} compares Cuhre with Suave for it, at a
target relative error of $10^{-3}$. The reference is Cuhre with the Filon rule
at 16 nodes and 32 panels and $\epsilon_{\text{rel}} = 10^{-5}$:
$F_L = 0.07582522(76)$, $F_T = 0.2611083(26)$.

\begin{table}[ht]
  \centering
  \small
  \begin{tabular}{lllll}
    \toprule
    & $F_L$ off by & $F_T$ off by & 1 core & 8 cores \\
    \midrule
    Cuhre, Filon (default)            & $0.0\sigma$ & $1.2\sigma$ & 0.8~s & 0.27~s \\
    Suave $n_{\text{new}} = 10^4$, Filon & $5$--$12\sigma$ & $6$--$7\sigma$ & 28--40~s & --- \\
    Suave $n_{\text{new}} = 10^5$, Filon & $0.5$--$1.2\sigma$ & $0.3$--$0.5\sigma$ & 23~s & 5.4~s \\
    Suave $n_{\text{new}} = 10^4$, by parts & up to $2.3\sigma$ & $3$--$4\sigma$ & 25~s & --- \\
    Suave $n_{\text{new}} = 10^5$, by parts & $0.4$--$1.8\sigma$ & $1.2$--$1.3\sigma$ & 13--15~s & 5.3~s \\
    \bottomrule
  \end{tabular}
  \caption{FE 5D with the large-$x$ extrapolation ($p = 4$), target relative
    error $10^{-3}$, two seeds for Suave. At $10^{-4}$, Cuhre took 4.0~s on one
    core and was within $0.1\sigma$ of the reference.}
  \label{tab:suave}
\end{table}

Suave is 15--50 times slower than Cuhre, and with its usual subdivision size
it is biased by several standard deviations without warning. Divonne is far too
slow. Cuhre with the Filon rule is therefore the recommended method.

On 8 cores, Cuhre gains a factor of 3 for the expensive Filon integrand. For the
cheap FE 4D and frozen-dipole integrands it is about 1.3 times faster on one
core than with 8 worker processes (\opt{cores} 0).

\section{Comparison with the Python reference}
\label{sec:python}

The Python references are vegas runs of the Python code with
$10^7$ points on 8 cores. For OT there is also a nested-quadrature value,
accurate to $10^{-10}$. For FE 5D the reference is the mean of three runs with
different seeds, with its error taken from their scatter. That scatter is about
twice the error vegas reports for one run, a consequence of the infinite
variance described in Section~\ref{sec:variables}. Python's FE results also
drift with the number of points: FE 4D $F_T$ is $0.38246(45)$ at $10^6$ points
and $0.38347(12)$ at $10^7$, $2.2\sigma$ apart.

\begin{table}[ht]
  \centering
  \small
  \begin{tabular}{lllll}
    \toprule
    & $F_L$ & pull & $F_T$ & pull \\
    \midrule
    \multicolumn{5}{l}{\textbf{OT}} \\
    Python, nested quadrature            & $0.0994504423$ & --- & $0.3913902122$ & --- \\
    Python, vegas $10^7$                 & $0.09945045(3)$ & & $0.3913905(2)$ & \\
    C++ Cuhre, $\epsilon_{\text{rel}} = 10^{-4}$ & $0.0994505(70)$ & & $0.391391(38)$ & \\
    C++ Cuhre, $\epsilon_{\text{rel}} = 10^{-7}$ & \multicolumn{4}{l}{agrees with the quadrature to $5\times10^{-9}$ relative} \\
    \midrule
    \multicolumn{5}{l}{\textbf{FE 4D}} \\
    Python, vegas $10^7$                 & $0.099269(44)$ & --- & $0.38347(12)$ & --- \\
    C++ Cuhre, $\epsilon_{\text{rel}} = 10^{-4}$ & $0.0993117(99)$ & $+0.9$ & $0.383456(38)$ & $-0.1$ \\
    C++ Vegas $10^6$, default variables   & $0.09931(18)$ & $+0.2$ & $0.38393(54)$ & $+0.8$ \\
    C++ Vegas $10^6$, Python variables    & $0.09910(32)$ & $-0.5$ & $0.3830(10)$ & $-0.5$ \\
    \midrule
    \multicolumn{5}{l}{\textbf{FE 5D}} \\
    Python, vegas, mean of three $10^7$ runs & $0.07606(47)$ & --- & $0.2769(11)$ & --- \\
    C++ Cuhre, $\epsilon_{\text{rel}} = 10^{-4}$ (by parts) & $0.0763908(64)$ & $+0.7$ & $0.276422(25)$ & $-0.4$ \\
    C++ Vegas $10^6$, default variables   & $0.07650(12)$ & $+0.9$ & $0.27705(38)$ & $+0.1$ \\
    C++ Vegas $10^6$, Python variables    & $0.0801(19)$ & $+2.0$ & $0.2875(43)$ & $+2.4$ \\
    \midrule
    \multicolumn{5}{l}{\textbf{Large $N_c$}, against Python vegas $10^6$} \\
    FE 4D, C++ Cuhre $10^{-3}$            & $0.099308(99)$ & $-0.1$ & $0.38352(38)$ & $-0.6$ \\
    FE 5D, C++ Cuhre $10^{-3}$            & $0.076406(64)$ & $+0.7$ & $0.27639(25)$ & $+0.6$ \\
    \bottomrule
  \end{tabular}
  \caption{Structure functions at $Q^2 = 2\ \text{GeV}^2$, $x_B = 10^{-3}$.
    The pull is the difference from the Python value over the combined
    error.}
  \label{tab:python}
\end{table}

Table~\ref{tab:python} shows that every C++ method agrees with Python within
its errors. The largest pulls, about $2\sigma$, are for FE 5D in the Python
variables, which share the infinite-variance problem of the Python code itself.
With the default settings all pulls are below $1\sigma$.

The C++ Monte Carlo errors are reliable: over 8 seeds the scatter of the
results was 0.6--1.5 times the reported error. From $10^6$ to $10^7$ points the
default-variable Vegas results moved by at most $2\sigma$, about $0.1\%$ for FE
5D, compared with the $2.2\sigma$ drift of Python's FE 4D.

\paragraph{The extrapolated dipole} has no Python reference. Three independent
methods agree for it: Cuhre with the Filon rule ($\epsilon_{\text{rel}} =
10^{-4}$: $F_L = 0.0758259(75)$, $F_T = 0.261111(26)$), Vegas by parts
($0.075761(82)$, $0.26119(23)$), and Vegas on the original 5D form in the Python
variables, i.e.\ the Python method ($0.07501(58)$, $0.2594(13)$). The pulls are
at most $1.4\sigma$. The extrapolation lowers FE 5D compared with the frozen
dipole: $F_L$ by 0.7\% and $F_T$ by 5.5\%.

\section{Performance}

Table~\ref{tab:speed} compares wall-clock times at matched precision. For the
C++ code, Cuhre's \opt{epsrel} was set to the relative error of the Python
result, and the time includes reading the BK file. The machine is an Apple M2
with 4 performance and 4 efficiency cores.

\begin{table}[ht]
  \centering
  \footnotesize
  \setlength{\tabcolsep}{3.5pt}
  \begin{tabular}{llllll}
    \toprule
    Precision of & \multicolumn{2}{l}{Python, 1 / 8 cores} & C++ & \multicolumn{2}{l}{speedup, 1 / 8 cores} \\
    \midrule
    OT, Python $10^6$ ($2\times10^{-6}$)      & 41.0~s & 10.6~s & 0.01~s & $\sim$4000 & $\sim$1000 \\
    FE 4D, Python $10^6$ ($1.2\times10^{-3}$) & 16.1~s & 9.1~s  & 0.86~s & 19 & 11 \\
    FE 5D, Python $10^6$ ($10^{-2}$)          & 25.8~s & 11.9~s & 0.24~s & 110 & 50 \\
    \textbf{all three}                        & 83~s   & 32~s   & 1.1~s  & \textbf{75} & \textbf{30} \\
    \midrule
    OT, Python $10^7$ ($3\times10^{-7}$)      & --- & 105~s & 0.01~s & --- & $\sim$10000 \\
    FE 4D, Python $10^7$ ($3\times10^{-4}$)   & --- & 71~s  & 2.6~s  & --- & 27 \\
    FE 5D, Python $10^7$ ($3\times10^{-3}$)   & --- & 95~s  & 0.62~s & --- & 150 \\
    \textbf{all three}                        & --- & 271~s & 3.3~s  & --- & \textbf{80} \\
    \bottomrule
  \end{tabular}
  \caption{Time to comparable precision (relative error in parentheses), frozen
    dipole. The C++ times are for one Cuhre process.}
  \label{tab:speed}
\end{table}

These numbers favour Python, because its reported FE errors are too small. For
many kinematic points, independent C++ runs in parallel give another factor of
about 4.7 on this machine: 8 concurrent FE 4D runs took 0.18~s per calculation
instead of 0.86~s. The gains come from compiled code (about 2--3 times per
point), Boost Bessel functions (1.8 times), computing L and T together (2 times),
the integration variables and the by-parts or Filon forms, and, above all for OT
and FE 5D, from replacing Monte Carlo by cubature.

\section{Recommendations and caveats}

\begin{itemize}
  \item The defaults (Cuhre, \opt{epsrel} $10^{-3}$, automatic FE 5D method)
    give about $10^{-3}$ relative precision in about a second per structure
    function. For final numbers use \opt{epsrel} $10^{-4}$, and check a point
    with a ten times smaller \opt{epsrel}.
  \item With the Filon rule below \opt{epsrel} $10^{-4}$, raise
    \opt{filon-nodes} and \opt{filon-panels}; at the defaults (8, 8) the
    interpolation error is far below the $10^{-4}$ integration error.
  \item All checks are at $x_B = 10^{-3}$. At smaller $x_B$ the FE kernel peak
    is narrower and $J_0$ oscillates more; spot-check a production point there
    against Vegas or a tighter \opt{epsrel}.
  \item Interpolating $T$ in $\xP$ above $x_0$ suits extrapolations written as
    functions of $x$. A different extrapolation that is not smooth in $x$ would
    need checking with more Filon nodes.
  \item For cross-checks, the Python method is reproduced by the options
    \opt{integrator vegas}, \opt{variables python} and
    \opt{fe5d-method original}.
\end{itemize}

\end{document}
